\subsection{La révélation par la pratique}

Mon intérêt pour la physique ne date pas de mon entrée en licence. 

Il est le fruit d’un long cheminement personnel, 
né d’une volonté de comprendre le monde qui m'entourait 
alors que j'exerçais une activité professionnelle qui ne me permettait pas d'exploiter ce potentiel. 

Pendant des années, en autodidacte, 
j'ai consacré mon temps libre à l'apprentissage des mathématiques et de la physique 
par la lecture, les cours en ligne et l'expérimentation personnelle. 

Ce qui était au départ une curiosité est devenu une évidence : 
je ne pouvais plus rester dans une pratique théorique isolée.

C'est cet engagement qui m'a conduit en licence de physique, 
et c'est ce même engagement que je retrouve aujourd'hui lors de mes travaux dirigés. 

Ici, la théorie rencontre enfin la réalité du terrain. 
J'ai compris que mon épanouissement professionnel se situe à l'intersection 
de ce savoir rigoureux acquis au fil des années 
et de ma capacité à \textbf{modéliser des systèmes pour les transformer}. 

Ce n'est pas une rupture avec mon passé, 
c'est la \textbf{concrétisation d'un projet} que je construis depuis longtemps.

\subsection{La valeur ajoutée de mon parcours}

Ma force ne réside pas dans une expérience professionnelle en physique, 
mais dans la synthèse entre mon parcours de terrain 
et ma passion pour les sciences. 

J’ai déjà été confronté à la réalité du monde professionnel ; 
j'y ai constaté la multitude de problèmes complexes 
qui attendent des solutions concrètes.

Cette expérience m'a appris une chose essentielle : 
la théorie, sans une finalité utile, reste isolée. 

Aujourd'hui, 
je souhaite articuler ma rigueur analytique 
— nourrie par des années de pratique personnelle en mathématiques et en physique — 
avec cette connaissance des rouages du terrain. 

Là où je me distingue, 
c'est dans cette volonté de \textbf{faire le lien}. 

Je ne cherche pas à appliquer la théorie pour la théorie, 
mais à \textbf{utiliser la puissance de modélisation} de la physique et des mathématiques 
pour résoudre les problèmes réels que j'ai pu observer. 

Mon ambition est simple : 
mettre ma capacité de réflexion au service de projets qui ont un \textbf{réel impact}. 

Pour un recruteur, 
cela signifie un collaborateur qui a déjà les pieds sur terre, 
une grande autonomie dans l'apprentissage, 
et une motivation profonde pour \textbf{l'utilité concrète} de ses actions.